\documentclass[11pt]{article}
\usepackage{tutorial}
\begin{document}
\tutorialtitle{1}{Temporal Point Processes and Races of Clocks}{Hazards, survival, likelihoods with silence, and why a race is a posterior over causes}

\section*{Why this tutorial}
Our models predict \emph{when} the next event happens and \emph{what} it is. The mathematical language for that is the
temporal point process. A \emph{race of clocks}, several candidate events competing to happen first, is the core primitive of
our model family. This tutorial builds the vocabulary from scratch: hazards, survival, the likelihood of an event sequence
(including the information carried by silence), marks, and races. It ends with the single most important identity for
learning: the gradient of a race's log-likelihood is ``responsibility minus exposure''.

\section{Events in continuous time}
A \emph{temporal point process} is a random sequence of event times $0<t_1<t_2<\dots$. Its counting process is
$N(t)=\#\{k: t_k\le t\}$. Write $\mathcal H_t$ for the history up to (but not including) $t$.

\begin{definition}[Conditional intensity / hazard]
The \emph{intensity} $\lambda(t)$ is the instantaneous event rate given the history:
\[
\lambda(t)\dd t = \Pr\{\text{an event in }[t,t+\dd t)\mid \mathcal H_t\}.
\]
For the waiting time $\tau$ since the last event, the same object is called the \emph{hazard} $h(\tau)$.
\end{definition}

The hazard determines everything else about the next waiting time:
\begin{align*}
\text{integrated hazard:}&\quad H(\tau)=\int_0^\tau h(s)\dd s,\\
\text{survival (no event yet):}&\quad S(\tau)=\Pr\{\tau_{\text{next}}>\tau\}=e^{-H(\tau)},\\
\text{density of the waiting time:}&\quad f(\tau)=h(\tau)\,S(\tau).
\end{align*}
\emph{Derivation.} $S(\tau+\dd\tau)=S(\tau)\,(1-h(\tau)\dd\tau)$, so $S'=-hS$ and $S=e^{-H}$; then $f=-S'=hS$.

\begin{example}[Exponential clock]
Constant hazard $h=\lambda$: $H=\lambda\tau$, $S=e^{-\lambda\tau}$, $f=\lambda e^{-\lambda\tau}$, mean $1/\lambda$.
The exponential is \emph{memoryless}: $\Pr\{\tau>a+b\mid \tau>a\}=e^{-\lambda b}$.
\end{example}
\begin{example}[Shaped clocks]
Weibull: $h(\tau)=\lambda k\tau^{k-1}$ (increasing hazard if $k>1$: ``due'' events). Log-normal: $\log\tau\sim\mathcal N(\mu,\sigma^2)$,
a \emph{delayed} clock that rarely fires early. Our models use mixtures of exponential and delayed (log-normal) clocks.
\end{example}

\section{The likelihood of a sequence, and the information in silence}
Observe events at $t_1,\dots,t_n$ in a window $[0,T]$. The probability density of exactly this observation is
\[
L=\Big(\prod_{k=1}^n \lambda(t_k)\Big)\exp\!\Big(-\int_0^T\lambda(t)\dd t\Big),\qquad
\log L=\sum_{k}\log\lambda(t_k)-\int_0^T\lambda(t)\dd t .
\]
The integral term is the log-probability of \emph{no events} in all the gaps, including the final stretch $(t_n,T]$.

\begin{keyidea}
Silence is data. The integral $\int\lambda$ penalises a model that expects events where none occurred. A model that ignores
silence can assign high likelihood to every observed event while predicting far too many unobserved ones. We call training
that keeps this term \emph{silence-aware supervision}.
\end{keyidea}

\paragraph{Marks.} If each event carries a type (mark) $k\in\{1,\dots,K\}$, give each mark its own intensity $\lambda_k(t)$.
The marked log-likelihood is $\sum_{\text{events}}\log\lambda_{k_j}(t_j)-\sum_k\int_0^T\lambda_k$. Equivalently, the total
rate $\Lambda=\sum_k\lambda_k$ generates times and the mark is drawn with probability $\lambda_k/\Lambda$.

\paragraph{Recording resolution.} Real timestamps are rounded (to seconds, days). A density model can put enormous density on the
rounding lattice and gain likelihood without predicting anything. Comparing models fairly requires scoring at a stated resolution
or \emph{dequantizing} (spreading each recorded time uniformly within its recording cell). This matters in practice: on some
public benchmarks a large part of published likelihood comes from the lattice.

\section{Races of clocks}
\begin{definition}[Race]
$K$ clocks run in parallel from the same start. Clock $i$ would fire after an independent waiting time with hazard $h_i$.
The race outputs the \emph{first} firing time $\tau=\min_i\tau_i$ and the \emph{winner} $I=\arg\min_i\tau_i$.
\end{definition}

\begin{proposition}[Race laws]
With $H_i(\tau)=\int_0^\tau h_i$:
\begin{enumerate}
\item The race's survival is the product $S(\tau)=\prod_i S_i(\tau)=e^{-\sum_i H_i(\tau)}$; its hazard is the sum $\sum_i h_i(\tau)$.
\item The joint density of (time, winner) is $f(\tau,i)=h_i(\tau)\,S(\tau)$.
\item Given that the first event happens at $\tau$, the winner is clock $i$ with probability
\[
\rho_i(\tau)=\frac{h_i(\tau)}{\sum_j h_j(\tau)} .
\]
\end{enumerate}
\end{proposition}
\begin{proof}
(1) Independence: no clock has fired by $\tau$ iff each has not. (2) Clock $i$ fires in $[\tau,\tau+\dd\tau)$ with probability
$h_i(\tau)S_i(\tau)\dd\tau$ and the others have not fired: multiply by $\prod_{j\neq i}S_j$. (3) Divide (2) by its sum over $i$.
\end{proof}

\begin{example}[Exponential race = softmax]
With constant rates $\lambda_i=e^{\theta_i}$, $\rho_i=e^{\theta_i}/\sum_j e^{\theta_j}$ does not depend on $\tau$, so the winner is a
softmax draw, independent of the time, and $\tau\sim\mathrm{Exp}(\sum\lambda_j)$. This is the ``Gumbel-max'' fact in disguise:
$-\log\tau_i=\theta_i+\text{Gumbel noise}$, and the arg-max of noisy scores is softmax-distributed.
\end{example}

\section{Learning: responsibility minus exposure}
Parametrize each clock's scale: $h_i(\tau)=e^{\theta_i}k_i(\tau)$, so $H_i=e^{\theta_i}K_i$ with $K_i=\int_0^\tau k_i$.
For one observed first event at $\tau$ (winner unobserved, or marginalized), $\log f(\tau)=\log\sum_j h_j(\tau)-\sum_j H_j(\tau)$. Then
\[
\boxed{\;\frac{\partial \log f(\tau)}{\partial\theta_i}=\rho_i(\tau)-H_i(\tau)\;}
\]
\emph{Responsibility} $\rho_i$ (the posterior probability that clock $i$ caused the event) minus \emph{exposure} $H_i$ (how much
clock $i$ was at risk without firing). If the winner is observed (a marked event), $\rho_i$ is replaced by the indicator $1[I=i]$.

\begin{keyidea}
The gradient of a race is an inference result: it attributes the observed event to its possible causes. That is literally
credit assignment, and at the output of a race it is exact and local: each clock needs only its own hazard and the race's total
rate. Note that the expected gradient is zero: $\E[\rho_i]=\E[H_i]$ (both equal the probability that $i$ wins), the usual
zero-mean property of a score.
\end{keyidea}

\section*{Exercises}
\begin{exercise} Show that the minimum of independent exponentials with rates $\lambda_i$ is exponential with rate $\sum\lambda_i$.\end{exercise}
\begin{solution} $\Pr\{\min>\tau\}=\prod e^{-\lambda_i\tau}=e^{-(\sum\lambda_i)\tau}$.\end{solution}
\begin{exercise} For two Weibull clocks with the same shape $k$, show the winner is independent of $\tau$. Is this true for different shapes?\end{exercise}
\begin{solution} $\rho_i=\lambda_i k\tau^{k-1}/\sum_j\lambda_j k\tau^{k-1}=\lambda_i/\sum\lambda_j$. With different shapes $\rho_i$ depends on $\tau$:
early events favour the clock with the decreasing hazard.\end{solution}
\begin{exercise} Write the log-likelihood of a sequence observed on $[0,T]$ with no events at all. What does maximizing it do?\end{exercise}
\begin{solution} $\log L=-\int_0^T\lambda$; it drives the predicted rate toward zero: silence teaches the model that events are rare.\end{solution}
\begin{exercise} Derive $\partial\log f/\partial\theta_i=\rho_i-H_i$ for the race above.\end{exercise}
\begin{solution} $\partial_{\theta_i}h_j=\delta_{ij}h_i$ and $\partial_{\theta_i}H_j=\delta_{ij}H_i$, so $\partial\log\sum h=h_i/\sum h=\rho_i$ and $\partial(-\sum H)=-H_i$.\end{solution}

\begin{inourmodels}
Race-of-clocks heads in \texttt{experiments/tpp/race\_tpp*.py} (EasyTPP battle B1, neural-TPP benchmark B4); silence-aware
supervision throughout; the recording-resolution floor of our clocks and the B4 lattice audit; theory note 160 \S2 (race credit is
posterior inference); note 168 (Fisher geometry of races, Tutorial 6).
\end{inourmodels}
\end{document}
